% !TeX root = ../tesis.tex

\chapter*{Introducción}
\addcontentsline{toc}{chapter}{\protect\numberline{}Introducción}	  		% Comment if you don't want the introduction to appear on the table of content. It will not have a number
\label{chapter:intro}

% this file is called up by thesis.tex
% content in this file will be fed into the main document

%
% --- JAUA ---
% Luz-materia puede confundir con propiedaes electrónicas de materiales, o de excitacines específicas (como plasmones o multipolos). Mejor habla del pronblema de espacimento por partículas.
% También omitamo biofotónica porque nunca usaste ese térmio en la tsis. Enuncia mejor diversas aplicaciones como imagenologhía, snesado, etc.
%

La interacción de la luz con la materia depende de propiedades como el tamaño, la geometría, la composición, la estructura interna y la orientación del sistema considerado \cite{bectondickinsonIntroductionFlowCytometry2000}. En este contexto, la biofotónica estudia el uso de la luz para analizar sistemas biológicos y desarrollar métodos de diagnóstico, monitoreo y tratamiento \cite{baldiniShiningLightFuture2025}. En el área de diagnóstico, la biofotónica aprovecha el hecho de que las alteraciones fisiológicas o patológicas en sistemas de células o en células individuales pueden modificar la manera en que interactúan con la luz, proporcionando información potencialmente útil para su caracterización \cite{baldiniShiningLightFuture2025}. %
% --- JAUA ---
% Corta acá el párrafo y comienza con el siguiente. Rescribe acorde
%

Un ejemplo de ello es el estudio del esparcimiento de luz en la sangre, cuya respuesta óptica puede mostrar cambios en las propiedades morfológicas y composicionales de sus componentes celulares asociados a distintas enfermedades \cite{ergulComputationalStudyScattering2010}.

Entre los principales componentes celulares de la sangre se encuentran los eritrocitos, cuya función fisiológica fundamental es el transporte de oxígeno \cite{palTextbookMedicalPhysiology2021}. En el intervalo %
% --- JAUA ---
% intervalo suena muy IA pocha. En el rango/espectro, o alfo .
% visibke, los eritrocitos, constituidos principalmente por hemoglobina, dominan la.....
%
 visible, estas células dominan la respuesta electromagnética de la sangre y están constituidas principalmente por hemoglobina \cite{bosschaartLiteratureReviewNovel2014}. En condiciones fisiológicas %
 % --- JAUA ---
 % Considerando el caso de una persona sana, los eritrocios...
 %
  presentan una geometría discoidal bicóncava; sin embargo,%
  % --- JAUA ---
  % innecesario el punto y coma, no??
  % cambia tamño por volumen
  distintas alteraciones hematológicas pueden modificar tanto su morfología como su composición \cite{palTextbookMedicalPhysiology2021}. Estas modificaciones pueden manifestarse, entre otros parámetros, como variaciones en la forma, el tamaño y la concentración de hemoglobina corpuscular media (CHCM) \cite{palTextbookMedicalPhysiology2021}.

%
% --- JAUA ---
% Mehir di que estas alteracion corresponden a respuestas espectrales y patrones de interferencia distintas, que podrían emplearse para caracterisar las enfermendas En particular, en el espectro visible, la hemoglobina cuenta con una estructa de bandas a longitudes de ondas caracterísicas que son relevantes para su fincón disiológica....
%
  El interés en estudiar estas modificaciones radica en que la respuesta electromagnética de los eritrocitos depende del efecto conjunto de su geometría y de su composición. 

En el intervalo visible, la hemoglobina presenta una estructura espectral característica, dominada por la banda de Soret y las bandas $\alpha$ y $\beta$, asociadas a transiciones electrónicas de la molécula \cite{dayerBandAssignmentHemoglobin2010}. Estas bandas producen variaciones en las partes real e imaginaria del índice de refracción y, por tanto, en los procesos de absorción y esparcimiento. Por ello, cambios en la morfología y en la CHCM pueden modificar tanto la respuesta espectral como la distribución angular de la luz esparcida por los eritrocitos \cite{bectondickinsonIntroductionFlowCytometry2000}. La sensibilidad de la respuesta óptica a estas propiedades ha motivado el desarrollo de técnicas ópticas para la identificación y diferenciación celular. Entre ellas, la citometría de flujo, que permite analizar grandes poblaciones de células a partir de la luz esparcida \cite{bectondickinsonIntroductionFlowCytometry2000}. Sin embargo, las señales registradas por un detector de citometría de flujo corresponden a la medición del campo esparcido dentro de una región angular determinada por la geometría del sistema de medición \cite{mckinnonFlowCytometryOverview2018}. En consecuencia, el conocimiento de la distribución angular completa del esparcimiento puede aportar información adicional para identificar regiones de observación particularmente sensibles a cambios en la morfología y en la composición de los eritrocitos.

La descripción teórica de esta respuesta constituye un problema electromagnético complejo %
% --- JAUA ---
% Es una ecuación diferencial que no oscila tonto en sus condicones de feontera, por qué sería complejo??  Hay que tener cuodado con los teŕminos que usamos,. 
%
 debido tanto a la geometría bicóncava de los eritrocitos como a que sus dimensiones son mayores que la longitud de onda de la radiación incidente. Por esta razón, el estudio del esparcimiento de luz por eritrocitos se ha abordado mediante distintos métodos numéricos. Lu et al. \cite{luSimulationsLightScattering2005}, por ejemplo, emplearon el método de \textit{Finite Difference Time Domain} (FDTD); Karlsson et al. \cite{karlssonNumericalSimulationsLight2005} emplearon el método de Discrete Dipole
Approximation (DDA) y Tsinopoulos et al. \cite{tsinopoulosLightScatteringAggregated2002} desarrollaron un enfoque basado en el Boundary Element Method (BEM). No obstante, los trabajos mencionados anteriormente se concentran principalmente en longitudes de onda cercanas a 632.8 nm, fuera de las principales bandas de absorción de la hemoglobina. Esto limita el estudio de la dependencia espectral de la respuesta óptica en regiones donde las propiedades ópticas de los eritrocitos presentan variaciones más pronunciadas. Por lo tanto, resulta de interés extender el análisis a un intervalo espectral que incluya la banda de Soret y las bandas $\beta$ y $\alpha$, considerando además las modificaciones morfológicas y composicionales asociadas a eritrocitos patológicos.

%
% --- JAUA ---
% Estos  cuatro parrafos estan bien escritos pero muy repetitvos.  Recuerda que acá no hay un resumen estricto de tu trabajjo entonees:
% Unidica los siguientes tres párrafos en uno solo: reduce el segundo párrafo porque no necesitamos información tan detallada. Y el pultimo párrafo reescielo omitiendo Capítulo 1, Capítulo 2.... Mejor pon sus títulos e hipervínculos a cada uno de ellos.

%
% --- JAUA ---
% De nuevo: todo esto lo tenemos que trabajar a profundidad de nuevo: te lo pongo muchas veces para que no creas que hemos acabado :)
%

A partir de esto, el objetivo de esta tesis es estudiar  cómo las modificaciones en la morfología y en la concentración de hemoglobina alteran las propiedades ópticas de los eritrocitos y %
% --- JAUA ---
% no pongas esto porque en principio y a lo hiciste -> si es posible 
identificar condiciones espectrales y angulares que aumenten el contraste entre células sanas y patológicas. Para ello, se estudian tres morfologías representativas: un eritrocito sano, un esferocito y un microcito. El eritrocito sano se modela mediante óvalos de Cassini, que permiten reproducir su forma bicóncava con un número reducido de parámetros geométricos; el esferocito se representa mediante una geometría esférica, mientras que el microcito conserva la forma bicóncava pero presenta dimensiones reducidas. A cada morfología se le asigna además una CHCM representativa, con el fin de incorporar las variaciones en su composición asociadas a estas alteraciones.

Para describir las propiedades ópticas de los eritrocitos se emplean datos experimentales del índice de refracción de la hemoglobina reportados en la literatura para distintas concentraciones. Debido a que estos datos se encuentran restringidos a ventanas espectrales finitas, se expanden las funciones dieléctricas al rango de 400 a 600 nm mediante un ajuste de osciladores de Lorentz para la parte imaginaria y la reconstrucción de la parte real mediante las relaciones de Kramers-Kronig sustractivas. Posteriormente, se estudia la convergencia numérica del método empleando como referencia la solución analítica de Mie para una esfera, para luego contrastar el modelo de eritrocito con resultados numéricos reportados en la literatura. Una vez establecidas las condiciones de simulación, se calculan las eficiencias y secciones transversales de extinción, absorción y esparcimiento, así como la sección transversal diferencial de esparcimiento y el factor de anisotropía en el intervalo espectral comprendido entre 400 y 600 nm. Finalmente, se analizan las bandas de Soret, $\beta$ y $\alpha$ y se estudia la distribución angular de la luz esparcida con el propósito de localizar intervalos de detección en los que aumente el contraste entre las distintas morfologías.

De esta manera, el objetivo general de esta tesis es caracterizar numéricamente las propiedades ópticas de eritrocitos sanos y con alteraciones morfológicas y determinar regiones espectrales y angulares potencialmente útiles para su discriminación. Como objetivos específicos se plantea construir modelos geométricos y ópticos representativos de las distintas células, verificar la convergencia y consistencia del método numérico, estudiar la dependencia espectral de las propiedades ópticas, comparar sus patrones angulares de esparcimiento y localizar ventanas de detección en las que las diferencias entre las morfologías sean más pronunciadas. 

Esta tesis está dividida en cuatro capítulos:  en el Capítulo 1 se presenta el formalismo electromagnético empleado para describir el esparcimiento de luz por partículas, incluyendo las ecuaciones de Maxwell, las secciones transversales y las relaciones de Kramers-Kronig y Kramers-Kronig sustractivas. En el Capítulo 2 se estudian la morfología, composición y principales alteraciones de los eritrocitos, así como la citometría de flujo y la construcción de las funciones dieléctricas empleadas en las simulaciones. En el Capítulo 3 se introduce el método de integración finita y se establece la configuración numérica mediante un análisis de convergencia. En el Capítulo 4 se presentan y discuten los resultados correspondientes al eritrocito sano, el esferocito y el microcito, incluyendo su respuesta espectral, sus patrones angulares de esparcimiento y la identificación de regiones de mayor contraste. Finalmente, se presentan las conclusiones generales y las principales líneas de trabajo a futuro.